\section{Introduction}
Fraud rings are groups of accounts that act together: they share devices,
payment cards or network addresses, and they move money among themselves in
cycles or fan-in and fan-out patterns. Graphs make these relations explicit,
which is why graph-based methods dominate fraud and anti-money laundering
(AML) research \cite{pourhabibi2020fraud,pan2025gbfrd}. Two problems appear as
soon as such methods run on a live stream of events.

\textbf{The hub dilemma.} Legitimate shared infrastructure, such as an office
gateway, a carrier-grade network address translation (NAT) address or a store
kiosk, links many unrelated accounts. Linking accounts by shared properties
produces sparse links through strongly identifying entities and over-linkage
through widely shared ones such as IP addresses \cite{yin2022behavioral}. The
usual remedy in rule-based systems is a hard degree cap: attributes shared by
more than $c$ accounts are ignored. The cap removes the noise but also opens an
evasion path, because a ring that shares a busy attribute becomes invisible to
it, much as fraudsters elsewhere hide by connecting to benign entities
\cite{liu2025robust}.

\textbf{Latency on a stream.} Accounts, links and transfers arrive
continuously. Re-running whole-graph queries on a timer delays detection, and
the cost of each pass grows with the graph rather than with the change, a
problem that incremental community detection addresses for clustering
\cite{zhang2026incremental,abughofa2021incremental}.

We make four contributions.
\begin{itemize}
\item \textbf{C1, hub-aware weighting.} We replace the degree cap with a
continuous weight: each shared attribute contributes $1/\ln(1+d)$, where $d$ is
its degree, with exponential time decay, and transfers inside the pair add
weight. The discount itself is FRAUDAR's column weighting
\cite{hooi2016fraudar}; our contribution is its use for shared-attribute ring
detection on a stream. Busy hubs count for little on their own, but a group that shares a hub
and also trades among itself still accumulates weight (Section~\ref{sec:method}).
What we can show is that keeping hub links is what matters: a constant weight
with its own threshold performs as well, and caps lose rings hidden behind hubs
(Section~\ref{sec:results}).
\item \textbf{C2, incremental detection.} A micro-batch engine re-evaluates only
the $k$-hop region touched by each batch, with fan-out limits that keep that
region bounded when it reaches a hub.
\item \textbf{C3, calibrated, explained alerts.} Candidate rings are scored by a
calibrated classifier, and each alert carries feature contributions, an
evidence subgraph and counterfactuals. The scorer supplies most of the
precision, for our candidates and for the baselines' (Section~\ref{sec:results}).
\item \textbf{C4, evaluation under evasion.} A deterministic simulator with
ring-level ground truth and five evasion strategies, including an attacker that
adapts to the detector's flags, lets us measure each component on identical
event streams (a paired design) with confidence intervals.
\end{itemize}

\section{Related Work}
\label{sec:related}
\textbf{Graph fraud detection.} Most recent graph fraud detection work trains
graph neural networks (GNNs) to label nodes. It addresses heterophily, where
fraudulent nodes connect mostly to benign ones \cite{hu2025hhgfd,lv2026mandate};
camouflage, by reweighting neighbours adversarially \cite{liu2025robust}; scarce
or unreliable labels, with one-class and semi-supervised training
\cite{li2025ocgfd,giles2023suspicious,peng2023csocan}; and interpretability
\cite{yu2024ingfd}. These methods output per-node scores and need labelled
training graphs. GB-FRD \cite{pan2025gbfrd} targets fraud rings in peer-to-peer
payment graphs, as we do, and Yin et al.\ \cite{yin2022behavioral} replace hard
shared-property links with behavioural similarity because of over-linkage
through shared entities. We keep shared-attribute links, which are cheap and
explainable, and change how much each one counts.

\textbf{Dense subgraphs and peeling.} Finding unusually dense groups is a
classical route to fraud and collusion detection: greedy peeling approximates
the densest subgraph \cite{charikar2000greedy}, spectral and greedy
densest-subgraph algorithms scale to billion-edge graphs \cite{feng2024unified},
and sequential peeling of a weighted bipartite graph exposes coordinated
adversarial annotators \cite{karaaslanli2025crowd}. FRAUDAR
\cite{hooi2016fraudar} makes peeling resistant to camouflage by weighting each
edge by a decreasing function of the degree of the object it points to,
recommending $1/\log(x+c)$ and arguing that $1/x$ shrinks too fast; Spade
\cite{jiang2022spade} maintains such peeling-based dense subgraphs
incrementally under edge insertions on transaction graphs. Core decomposition,
the peeling primitive we use to extract dense cores, has efficient parallel and
GPU implementations \cite{ahmad2023gpukcore,li2021graphblas} and is related to
densest-subgraph approximation \cite{dhulipala2022dp}. Our degree discount is
therefore not new: it is FRAUDAR's column weighting, and close to
Adamic--Adar. What we study is applying it to shared-attribute ring detection
on a stream, as a replacement for the degree cap of rule-based systems, with
incremental detection, ring tracking and calibrated alerts around it.

\textbf{Community detection on dynamic graphs.} Louvain modularity optimisation
\cite{blondel2008louvain} remains a standard way to find groups and is still
being refined \cite{li2023dplouvain}; we use it as a baseline. Incremental community
detection avoids repartitioning a changing graph by updating only the
affected communities \cite{zhang2026incremental,abughofa2021incremental}.
Community structure, however, is not ring structure: a benign hub joins
thousands of users into one community.

\textbf{Similarity weighting.} Local similarity indices such as Jaccard and
Adamic--Adar \cite{adamic2003friends}, which discounts a shared neighbour by
the logarithm of its degree, are standard in link prediction
\cite{varma2022jaccard}. Our hub weight applies a similar discount to shared
attributes, and we match flagged rings to true ones by Jaccard overlap.

\textbf{AML on public data.} Work on the Elliptic Bitcoin graph compares tabular
ensembles with GNNs under temporal splits
\cite{badjie2026ensembles,kumar2026tagat,lu2026lagcn,shah2025aml}; AMLworld
provides synthetic laundering attempts with typology labels
\cite{altman2023amlworld}. These studies classify transactions, so their
scores are not comparable with our ring-level metrics
(Section~\ref{sec:public}).

\textbf{Calibration and explanations.} Post-hoc calibration by isotonic
regression \cite{zadrozny2002transforming} and its comparison with Platt
scaling \cite{niculescu2005predicting} make classifier scores usable as
probabilities; Azbergen et al.\ \cite{azbergen2026calibration} compare such
calibrators for AML risk scores, and learned calibration has been studied
beyond classification \cite{ding2021temperature}. Counterfactual explanations
state what would have to change for a different decision; their robustness
\cite{artelt2021robustness}, consistency across model updates
\cite{movin2025acdc} and per-feature importance \cite{meulemeester2026cfi} are
active topics. We attach simple, model-specific counterfactuals to every alert.

\section{Method}
\label{sec:method}
Fig.~\ref{fig:arch} shows the system. A seeded simulator emits observable
events (account joins and leaves, attribute links, transfers); ground truth
goes only to a ledger the detector never sees.

\begin{figure*}[t]
\centering
\includegraphics[width=\textwidth]{\figpath/architecture.pdf}
\caption{Overview of the system. (a) A seeded simulator emits the observable
event stream in micro-batches; labels go only to a ledger. (b) Each batch is
written to the graph store and expanded to the bounded region it touched.
(c) Shared attributes are weighted by their degree, so a busy hub counts for
little and a rarely shared device for much. (d) Incremental detectors run on
the region and a tracker confirms rings across batches. (e) Confirmed rings are
scored, explained and alerted. Solid arrows: data flow; dashed arrows: labels
(evaluation only) and, in closed-loop runs, alerts returned to the attacker.}
\label{fig:arch}
\end{figure*}

\subsection{Graph model}
Users link to devices, cards and IP addresses
(\texttt{USED\_DEVICE}, \texttt{HAS\_CARD}, \texttt{LOGGED\_IN\_FROM}) and send
transfers to each other. The same contract is implemented over Neo4j (Cypher)
and over an in-memory store with identical semantics; all experiments below use
the in-memory store. In live runs, events reach the engine through an
in-process queue or Redis Streams, whose consumer settings (batch size,
blocking timeout) trade throughput against latency \cite{bora2026hyades}.

\subsection{Hub-aware weighting (C1)}
For two users $u,v$ the edge weight is
\begin{equation}
\begin{split}
W(u,v)={}&\sum_{a\in A(u)\cap A(v)} h(\deg a)\,\delta(t_a)\\
&+\min\Big(\beta\sum_{x\in T(u,v)}\delta(t_x),\,\beta_{\max}\Big),
\end{split}
\label{eq:w}
\end{equation}
where $h(d)=1/\ln(1+d)$, $\delta(t)=2^{-t/\tau}$ is exponential decay with
half-life $\tau=300$\,s, $t_a$ is the age of the older of the two links to
attribute $a$, and $T(u,v)$ are the transfers between $u$ and $v$ in the
window, with $\beta=0.15$ and $\beta_{\max}=0.6$. A device shared by three
accounts weighs $0.72$; an office address shared by fifty weighs $0.25$.
Candidate rings are the connected components of the graph of pairs with
$W\ge\theta_e$ (and a non-zero attribute term), peeled to cores of edge density
at least $0.75$ and size 3--25. The edge threshold $\theta_e$ is selected on
validation seeds (Section~\ref{sec:protocol}).

\subsection{Incremental detection (C2)}
Events are grouped into micro-batches of 50. For each batch the engine takes
the users it touched, expands them by $k=2$ hops through shared attributes and
transfers, and builds a view of that region over a 900\,s window, together with
the \emph{global} degree of every attribute in it, so $h(\deg a)$ is the same as
in a full pass. The first hop may cross attributes with up to 250 users (so the
co-members of a hub are re-examined); later hops cross only attributes with at
most 20 users, which stops the region from spreading through hubs. The first-hop
limit is itself a degree cap on the region: a pair whose only link is an
attribute shared by more than 250 users never enters the same region. For the
hub-aware weight this costs nothing, since one such attribute contributes at most
$1/\ln 252\approx0.18$, below the edge threshold, so a full pass would not link
the pair either; a weight that lets one busy attribute reach the threshold, such
as a constant, does lose such pairs (Section~\ref{sec:results}). Incremental
detectors (hub-aware shared attribute, temporal cycle, smurfing) run on the
view; whole-graph detectors (the static Cypher rules and community detection)
run on a timer.

\subsection{Ensemble, tracking and scoring (C3)}
Candidates from different detectors are merged when their Jaccard overlap is at
least $0.5$, or when one contains $80\%$ of the other and they differ in size by
at most a factor of two. A tracker matches candidates to existing rings across
batches; a ring is \emph{confirmed} after two consecutive hits and dissolved
after three misses. Confirmed rings are scored by a logistic-regression model
over 19 ring features (size, density, which attribute kinds are shared, the
hub-adjusted weight, cycle presence and length, burstiness, account ages,
internal transfers per member, amount statistics, growth, hub exposure and
the number of agreeing detectors), calibrated with isotonic regression
\cite{zadrozny2002transforming}. A ring becomes an alert when its calibrated risk
reaches the alert threshold $\theta_a$. Each alert carries the per-feature
contributions to the score, the subgraph of shared attributes and transfers
that formed it, and counterfactuals (for the leading features, the value at which
the risk would fall below the threshold, and for hub-aware rings the degree the
shared attribute would need before the rule stops firing). The model is trained on
seeds disjoint from all evaluation seeds (Fig.~\ref{fig:training}): the detector
runs on six scenarios with seeds 101--110, each confirmed or changed ring
becomes a row of its 19 features, labelled from the hidden ledger, and within
each run the first 70\% of rows fit the regression and the last 30\% the
calibrator.

\begin{figure*}[t]
\centering
\includegraphics[width=\textwidth]{\figpath/training_pipeline.pdf}
\caption{Training and use of the ring scorer. (a)--(b) The simulator and
detector run on training seeds, and ring snapshots are labelled from the
ground-truth ledger, which the detector never reads. (c) Each snapshot is a row
of 19 ring features. (d) Within each run, rows in time order are split 70/30
between logistic regression and isotonic calibration. (e) Held-out test seeds
(values from \texttt{models/ring\_scorer\_lr.json}). Below: in detection, each
changed ring is scored every batch, and an alert carries its explanation;
analyst feedback adjusts the threshold.}
\label{fig:training}
\end{figure*}

\section{Experimental Setup}
\label{sec:setup}
\subsection{Simulator and scenarios}
The simulator is a deterministic discrete-event generator: a scenario file and
a seed fix the whole stream, which we verify by a digest of the event sequence
recorded with every result. Benign users join and leave, link to devices, cards
and IPs drawn from heavy-tailed popularity, and transact; benign hubs (offices,
carrier NAT, kiosks, family devices) are generated with configured sizes. Ring
types are shared device, card or IP, mixed attribute, money cycle, layered
chain and smurfing fan-in/fan-out. Evasion strategies are hub hiding (the ring
shares a busy benign hub), low-and-slow activity, identity rotation, decoy edges
to benign users, and an adaptive attacker that shifts its strategy mix away
from strategies that were flagged.
Fig.~\ref{fig:simulator} shows the generator and its two outputs, and
Table~\ref{tab:scenario-def} lists the scenarios.

\begin{figure*}[t]
\centering
\includegraphics[width=\textwidth]{\figpath/simulator_design.pdf}
\caption{Simulator design. (a) A scenario file and a seed fix the whole run.
(b) One seeded generator combines a benign population, benign hubs, planted
rings and an evasion layer. (c) Observable events go to the detector, with a
digest of the sequence recorded with every result; ring lifecycle events go
only to the ground-truth ledger used for evaluation and the Oracle view. In
closed-loop runs, flags are returned to the adaptive attacker.}
\label{fig:simulator}
\end{figure*}

\begin{table}[t]
\centering
\caption{Scenarios (from \texttt{configs/scenarios/*.yaml}).}
\label{tab:scenario-def}
\small
\setlength{\tabcolsep}{3pt}
\begin{tabular}{lrrrl}
\toprule
Scenario & Users & Dur.\,(s) & Ev./s & Rings and evasion \\
\midrule
baseline & 30 & 60 & 10 & 4 planted, none \\
high\_noise & 150 & 60 & 50 & 4 planted + random, decoys \\
hub\_heavy & 100 & 60 & 25 & 4 planted, hub hiding \\
adversarial & 120 & 90 & 30 & 5 planted + random, all five \\
drift & 200 & 90 & 40 & 2 planted + random, hub hiding \\
demo\_story & 120 & 90 & 20 & 3 planted, hub hiding \\
scale\_sweep & $10^3$+ & 120 & 100 & latency benchmark (E4) \\
\bottomrule
\end{tabular}
\end{table}

\subsection{Protocol}
\label{sec:protocol}
Seeds are split three ways: scorer training (101--110, test 201--205),
threshold selection (validation, 301--303) and evaluation (test, 1001 and up).
On the validation seeds we evaluated a grid of alert thresholds
$\theta_a\in\{0.2,0.3,0.5,0.7,0.8,0.9\}$ and edge thresholds
$\theta_e\in\{0.2,0.3,0.4,0.5,0.6\}$ on the five accuracy scenarios; the grid was
widened whenever the best value sat on its edge. The selected values are
$\theta_a=\TunedAlert$ and $\theta_e=\TunedEdge$ (mean validation F1
$\TunedValF$). We report the proposed method both with these values and with
the a-priori values $\theta_a=\theta_e=0.5$.

Every method in an experiment processes the \emph{identical} event stream for a
given scenario and seed (the attacker does not receive feedback in these runs),
so differences between methods are paired. We report the mean and 95\%
$t$-interval over runs, and test differences with the Wilcoxon signed-rank test
on paired runs, with Cohen's $d_z$ as effect size.

\subsection{Metrics}
A flagged ring \emph{matches} a true ring when their member sets overlap with
Jaccard at least $0.5$. Ring precision is the fraction of flagged rings that
match a true ring; ring recall is the fraction of true rings matched at some
point while active. A run in which nothing is flagged has undefined precision;
we count its F1 as 0. Time to detect (TTD) is the time from a ring's
activation to its first matching flag. Detection latency is the wall-clock time
per micro-batch for the dirty-region expansion, the view and the incremental
detectors.

\subsection{Baselines}
\begin{itemize}
\item \textbf{Static Cypher, cap 100 / cap 10.} The original rules: shared
device (Q1), shared card (Q2), 3-cycles of transfers (Q3), and shared IP with a
degree cap (Q4a/Q4b), run on the whole graph every 5\,s of simulated time.
\item \textbf{Louvain.} Louvain communities \cite{blondel2008louvain} of the
user graph projected through shared attributes and transfers, kept when dense
enough.
\item \textbf{Connected components.} Weakly connected components of the same
projection, no density filter.
\end{itemize}
All baselines use the same two-hit confirmation as the proposed method.

\subsection{Environment}
All runs used one laptop (Intel Core i3-1115G4, 2 cores / 4 threads, 7.7\,GB
RAM, Windows 11), Python 3.11.9, NumPy 2.2.6, scikit-learn 1.7.2,
NetworkX 3.5, with experiments run one at a time.

\section{Results}
\label{sec:results}

\subsection{RQ1: accuracy against baselines}
% generated from results/e1_baselines/results.csv (sha256 of LF text 898355fc90852cc6, 300 rows)
% generated from results/e1b_scored_baselines/results.csv (sha256 of LF text 0d8b65d8aef12277, 300 rows)
\begin{table*}[t]
\centering
\caption{Ring-level detection on 5 scenarios $\times$ 10 test seeds, pooled (mean $\pm$ 95\% CI over runs;
a run with no flags has undefined precision and counts as F1 $=0$). TTD: median time from ring activation to
first flag. Latency: median per-batch detection time on the in-memory store; $^\dagger$full-graph
methods run on a 5\,s timer, not per batch, so their entry is the median time of one full-graph pass. $p$: paired Wilcoxon signed-rank
test of F1 against the proposed method on identical streams. Scored baselines use the shipped scorer (\texttt{models/ring\_scorer\_lr.json}), with the alert threshold selected per method on the validation seeds; the proposed method includes the same scorer, so only rows within a group isolate the candidate generator. $p$ values are Holm-corrected.}
\label{tab:baselines}
\begin{tabular}{lcccccc}
\toprule
Method & Precision & Recall & F1 & TTD (s) & Latency (ms) & $p$ \\
\midrule
\multicolumn{7}{l}{\textit{Without a scorer}} \\
Static Cypher (cap 100) & 0.12 $\pm$ 0.03 & 0.61 $\pm$ 0.08 & 0.19 $\pm$ 0.05 & 9.5 & 6.0$^\dagger$ & $<$0.001 \\
Static Cypher (cap 10) & 0.14 $\pm$ 0.04 & 0.61 $\pm$ 0.08 & 0.21 $\pm$ 0.05 & 9.5 & 5.9$^\dagger$ & $<$0.001 \\
Louvain communities & 0.39 $\pm$ 0.09 & 0.33 $\pm$ 0.08 & 0.29 $\pm$ 0.06 & 18.8 & 38.4$^\dagger$ & $<$0.001 \\
Connected components & -- & 0.00 $\pm$ 0.00 & 0.00 $\pm$ 0.00 & -- & 11.1$^\dagger$ & $<$0.001 \\
\midrule
\multicolumn{7}{l}{\textit{With the calibrated scorer}} \\
Static Cypher (cap 100) + scorer & 0.63 $\pm$ 0.08 & 0.48 $\pm$ 0.09 & 0.52 $\pm$ 0.08 & 8.6 & 4.5$^\dagger$ & $<$0.001 \\
Static Cypher (cap 10) + scorer & 0.64 $\pm$ 0.07 & 0.53 $\pm$ 0.09 & 0.56 $\pm$ 0.07 & 9.3 & 4.6$^\dagger$ & $<$0.001 \\
Louvain communities + scorer & 0.81 $\pm$ 0.12 & 0.32 $\pm$ 0.09 & 0.41 $\pm$ 0.09 & 19.0 & 27.9$^\dagger$ & $<$0.001 \\
Hub-aware incremental (proposed) & 0.76 $\pm$ 0.06 & 0.74 $\pm$ 0.06 & \textbf{0.73 $\pm$ 0.05} & 9.3 & 26.6 & -- \\
Proposed, a-priori thresholds & 0.73 $\pm$ 0.06 & 0.69 $\pm$ 0.06 & 0.68 $\pm$ 0.05 & 9.0 & 22.7 & 0.034 \\
\bottomrule
\end{tabular}
\end{table*}

\input{tables/tab1b_scenarios.tex}

Table~\ref{tab:baselines} pools \EScenarios{} scenarios $\times$ \ESeeds{}
test seeds for every method, all on identical streams. Without a scorer, the
static rules reach ring F1 \ECapTenF{} with a degree cap of 10 and
\ECapHundredF{} with a cap of 100, Louvain \ELouvainF{} and connected
components \EPlainCCF{}. These rows differ from the proposed method in two
ways at once: how candidates are generated and whether they are scored. With
the same calibrated scorer and an alert threshold selected per method on the
validation seeds, the static rules reach \ECapTenSF{} (cap 10) and
\ECapHundredSF{} (cap 100) and Louvain \ELouvainSF{}, against \EPropF{}
($\pm$\EPropFCI{}) for the proposed method; every difference to the proposed
method remains significant after Holm correction ($p<0.001$). Within the
scored group the differences isolate the candidate generator as a whole
(Section~\ref{sec:attribution} separates its parts); between a generator's
unscored and scored rows they isolate the scorer. The shipped scorer was
trained on the proposed pipeline's candidates, which favours it;
Section~\ref{sec:existing} gives every method its own scorer. For the rules with cap 10, the scorer raises F1
from \ECapTenF{} to \ECapTenSF{}, through precision (\ECapTenP{} to
\ECapTenSP{}) at some cost in recall (\ECapTenR{} to \ECapTenSR{}); the
candidate generator raises it further to \EPropF{}, mostly through recall
(\ECapTenSR{} to \EPropR{}). The scorer was trained on candidates of the proposed pipeline, so
on rule and Louvain candidates it works under some distribution shift. The
unscored static rules reach recall
\ECapTenR{} but precision only \ECapTenP{}: the shared-attribute rules and the
3-cycle rule fire on many benign groups, and two-hit confirmation does not
remove them. Connected components never produce a group within the ring size
limits, because benign hubs join almost all users into one component; Louvain
splits that component, but its communities rarely coincide with a ring. The
thresholds selected on validation seeds improve on the a-priori ones (F1
\EPropF{} vs.\ \EPropDefaultF{}, $p=0.034$, Holm over the eight comparisons
of Table~\ref{tab:baselines}).

Per scenario (Table~\ref{tab:scenarios}), the gap to the static rules is
largest on \texttt{drift} and \texttt{high\_noise}, where many benign groups
satisfy the rules and their precision collapses, and smallest on the clean
\texttt{baseline} scenario and on \texttt{adversarial}, which is the hardest
scenario for every method. Against the scored rules the difference depends on
the scenario: \RfOldNoiseD{} on \texttt{high\_noise}, but only \RfOldHubD{}
(95\% CI \RfOldHubLo{} to \RfOldHubHi{}, wins/ties/losses \RfOldHubWTL{}) on
\texttt{hub\_heavy}, the scenario built for hub hiding. With their own scorer
the rules reach \RfScRulesTabF{} on these streams.

\textbf{Degree cap vs.\ hub weighting.} Fig.~\ref{fig:frontier} isolates the
mechanism this paper replaces. Only shared-attribute and mixed rings count as
ground truth here (flags that match only other ring types are ignored), the
static rules run without the cycle rule, and the hub-aware detector runs alone.
Without the scorer, the rules' recall is \FrCapMaxR{} for every cap from 5 to
200, although caps of 50 and more exceed every benign hub in these scenarios: the
rules miss these rings because each rule returns all users of one attribute,
so a ring inside a larger group fails the Jaccard match, whatever the cap.
Section~\ref{sec:attribution} tests the cap where it does act, inside the
pairwise detector. Hub weighting reaches
recall \FrHubMaxR{}, but on its own at precision \FrHubPAtMaxR{}. With the
calibrated scorer on top, hub weighting reaches recall \FrHubSMaxR{} (95\%
bootstrap interval \RfFrHubSRLo{}--\RfFrHubSRHi{}; precision \FrHubSPAtMaxR{}) and its best F1 is \FrHubSBestF{} (alert threshold
\FrHubSBestFAt{}); the cap with the scorer stays below recall \FrCapSMaxR{}
(\RfFrCapSRLo{}--\RfFrCapSRHi{}) and
its best F1 is \FrCapSBestF{}. The difference in best F1, \FrDiff{}, has a
95\% bootstrap interval of \FrDiffLo{} to \FrDiffHi{} over the 20 runs (the best
threshold re-selected in every resample; whole runs resampled within scenario),
so at the best operating point the two mechanisms are not distinguishable on
F1. The best operating point is chosen on these test runs for both arms, so
both best-F1 values are optimistic. At low recall the cap with the scorer is the
more precise of the two (\FrCapSBestFP{} at recall \FrCapSBestFR{}). Hub
weighting therefore does not dominate the cap everywhere; it extends the
frontier to recall levels the cap cannot reach, and the scorer is what makes
those levels usable.

\begin{figure}[t]
\centering
\includegraphics[width=\columnwidth]{\figpath/fig1_pr_frontier.pdf}
\caption{Precision and recall on shared-attribute rings (micro-averaged over
4 scenarios $\times$ 5 test seeds). Solid: detector alone, swept over the
degree cap or the edge threshold $\theta_e$. Dashed: with the calibrated scorer,
swept over the alert threshold $\theta_a$ (cap 10; $\theta_e=\TunedEdge$).
Best-F1 difference with the scorer (hub minus cap), each at the best alert
threshold on these test runs (an oracle choice): \FrDiff{}, 95\% bootstrap
interval \FrDiffLo{} to \FrDiffHi{} (\texttt{results/e3b\_frontier\_ci}).}
\label{fig:frontier}
\end{figure}

\subsection{Comparison with existing methods}
\label{sec:existing}
% generated by scripts/make_refine_tables.py from results/w2_eval, results/w3b_eval
\begin{table*}[t]
\centering
\caption{Comparison with existing methods on fresh test seeds 2101--2120 (5 scenarios $\times$ 20
seeds, paired). Every method has a scorer trained on its own candidates (seeds 101--110) unless
noted, and its structural parameter and alert threshold selected on seeds 301--303 with at least
the proposed method's budget. PR AUC: ring-level, threshold-free (not defined without a ring
score). $\Delta$: proposed minus row (scenario-weighted, 95\% cluster bootstrap), $p$:
stratified sign-flip; the comparison with the strongest baseline by validation F1 (GraphSAGE, no
scorer) is confirmatory, the others exploratory. $^\ddagger$Exploratory variant added after the
degree-cap control. $^\S$Selection uses the scenario's identity (an oracle advantage).}
\label{tab:existing}
\small
\setlength{\tabcolsep}{2.5pt}
\resizebox{\textwidth}{!}{%
\begin{tabular}{lcccccccccc}
\toprule
Method & P & R & F1 & PR AUC & Hub-hiding & Decoy & Adaptive & $\Delta$ & 95\% CI & $p$ \\
\midrule
Proposed (shipped) & 0.74 & 0.72 & 0.719 $\pm$ 0.040 & 0.780 & 65/105 & 50/59 & 24/28 & -- &  &  \\
Proposed, constant weight$^\ddagger$ & 0.81 & 0.71 & 0.744 $\pm$ 0.043 & 0.755 & 69/105 & 50/59 & 11/28 & $-$0.025 & [$-$0.054, +0.003] & 0.192 \\
\midrule
GraphSAGE, no ring scorer & 0.88 & 0.72 & 0.778 $\pm$ 0.040 & -- & 73/105 & 37/59 & 18/28 & $-$0.059 & [$-$0.097, $-$0.023] & 0.025 \\
GraphSAGE + scorer & 0.85 & 0.69 & 0.751 $\pm$ 0.041 & 0.719 & 69/105 & 37/59 & 16/28 & $-$0.033 & [$-$0.072, +0.007] & 0.227 \\
Rules, cap tuned per scenario$^\S$ & 0.87 & 0.58 & 0.659 $\pm$ 0.053 & 0.619 & 71/105 & 23/59 & 22/28 & +0.059 & [+0.023, +0.094] & 0.014 \\
Rules, cap and cycle rule tuned & 0.91 & 0.55 & 0.646 $\pm$ 0.055 & 0.655 & 54/105 & 23/59 & 18/28 & +0.072 & [+0.037, +0.106] & 0.003 \\
Louvain, resolution tuned & 0.61 & 0.48 & 0.521 $\pm$ 0.051 & 0.562 & 43/105 & 28/59 & 11/28 & +0.198 & [+0.158, +0.239] & $<$0.001 \\
Leiden, resolution tuned & 0.86 & 0.40 & 0.507 $\pm$ 0.054 & 0.678 & 5/105 & 4/59 & 4/28 & +0.212 & [+0.172, +0.252] & $<$0.001 \\
FRAUDAR-style peeling (= Spade output) & 0.85 & 0.37 & 0.456 $\pm$ 0.058 & 0.519 & 58/105 & 2/59 & 5/28 & +0.263 & [+0.221, +0.304] & $<$0.001 \\
Gradient boosting on Louvain groups & 0.91 & 0.33 & 0.427 $\pm$ 0.062 & 0.510 & 0/105 & 7/59 & 3/28 & +0.291 & [+0.249, +0.334] & $<$0.001 \\
\midrule
\multicolumn{11}{l}{\textit{Hybrid, fresh seeds 2301--2320 ($\Delta$: hybrid minus row)}} \\
Proposed + GraphSAGE detector & 0.76 & 0.80 & 0.770 $\pm$ 0.036 & 0.808 & 89/107 & 59/67 & 23/30 & -- &  &  \\
GraphSAGE, no ring scorer & 0.88 & 0.72 & 0.780 $\pm$ 0.037 & -- & 68/107 & 50/67 & 16/30 & $-$0.010 & [$-$0.044, +0.023] & 0.664 \\
GraphSAGE + scorer & 0.84 & 0.70 & 0.755 $\pm$ 0.039 & 0.701 & 69/107 & 47/67 & 14/30 & +0.015 & [$-$0.021, +0.051] & 0.547 \\
Proposed (shipped) & 0.76 & 0.72 & 0.718 $\pm$ 0.037 & 0.792 & 68/107 & 54/67 & 20/30 & +0.052 & [+0.023, +0.081] & 0.003 \\
\bottomrule
\end{tabular}}
\end{table*}

Table~\ref{tab:existing} compares the pipeline with six existing methods under
the protocol of Section~\ref{sec:setup}, on test seeds no experiment had used.
Against every non-learned method it is ahead: by \RfExrulesperscenarioDA{} ring
F1 (95\% CI \RfExrulesperscenarioLo{} to \RfExrulesperscenarioHi{}) against the
rules with a cap and cycle rule tuned per scenario, which is the strongest rule
baseline and uses the scenario's identity, by \RfExrulesDA{} against the rules
tuned globally, and by \RfExlouvainDA{}, \RfExleidenDA{} and \RfExpeelingDA{}
against Louvain, Leiden and FRAUDAR-style peeling. With their own scorers the
rules come much closer than in Table~\ref{tab:baselines} (F1 \RfExrulesF{}
against \ECapTenSF{} with the shipped scorer). Peeling and Leiden are precise
but find few rings, and peeling and Leiden find almost no decoy rings
(\RfExpeelingDecoy{} and \RfExleidenDecoy{}).

A GraphSAGE classifier trained on ring-membership labels is the one method
ahead of the pipeline: F1 \RfExgnnunscoredF{} against \RfExproposedF{}
(difference \RfExgnnunscoredD{}, 95\% CI \RfExgnnunscoredLo{} to
\RfExgnnunscoredHi{}, confirmatory). With a ring scorer of its own GraphSAGE
reaches \RfExgnnF{}, and the difference (\RfExgnnD{}) is not significant. The
gap is not uniform: GraphSAGE is ahead on \texttt{drift} (\RfExgnnunscoredDrift{})
and \texttt{high\_noise} (\RfExgnnunscoredNoise{}), behind on
\texttt{adversarial} (\RfExgnnunscoredAdv{}), and the pipeline finds more
adaptive rings (\RfExproposedAdapt{} against \RfExgnnunscoredAdapt{}) and decoy
rings (\RfExproposedDecoy{} against \RfExgnnunscoredDecoy{}). Ranked by their
calibrated scores, the pipeline's rings have the highest ring-level PR AUC of
all scored methods (\RfExproposedPR{} against \RfExgnnPR{} for GraphSAGE with a
scorer; difference \RtPrD{}, 95\% CI \RtPrLo{} to \RtPrHi{}, exploratory), so
GraphSAGE's advantage lies in its operating point, not in ranking.
All of these data come from one generator, and GraphSAGE learns its regularities
from about 320{,}000 labelled user snapshots.

\textbf{Hybrid.} Adding the GraphSAGE classifier as a fourth detector inside the
pipeline, with a ring scorer trained on the combined candidates, was decided
after this comparison, so we tested it on another fresh block (2301--2320,
bottom of Table~\ref{tab:existing}). The hybrid reaches F1 \RfHyhybridF{}
against \RfHygnnunscoredF{} for GraphSAGE alone (hybrid minus GraphSAGE
\RfHygnnunscoredD{}, 95\% CI \RfHygnnunscoredLo{} to \RfHygnnunscoredHi{}) and
\RfHyproposedF{} for the pipeline alone. Any gain of the hybrid belongs to the
combination, not to a single element of the pipeline.

\subsection{What the advantage depends on}
\label{sec:attribution}
% generated by scripts/make_refine_tables.py from results/r1_cap_control, results/nd_evidence_count
\begin{table*}[t]
\centering
\caption{What the pairwise detector's advantage depends on: isolation controls inside the full
pipeline (same detectors, tracker and scorer; 5 scenarios $\times$ 10 or 20 fresh test seeds).
$\Delta$: scenario-weighted mean paired difference in ring F1 (A minus B) with 95\% cluster
bootstrap interval; $p$: stratified sign-flip test; (H): confirmatory comparison (Holm family of
four); others exploratory. Last column: recall of hub-hiding rings, A / B.}
\label{tab:attribution}
\small
\setlength{\tabcolsep}{3pt}
\resizebox{\textwidth}{!}{%
\begin{tabular}{llcccccc}
\toprule
A vs.\ B & Seeds & F1 A / B & $\Delta$ & 95\% CI & $p$ & $\Delta$ hub\_heavy & Hub-hiding \\
\midrule
Keep hub links vs.\ cap 10, $1/\ln(1+d)$ (H) & 2001--10 & 0.700 / 0.616 & +0.084 & [+0.048, +0.121] & 0.002 & +0.34 & 29/52 / 9/52 \\
Keep hub links vs.\ cap 10, constant & 2001--10 & 0.729 / 0.671 & +0.057 & [+0.039, +0.077] & 0.007 & +0.31 & 31/52 / 7/52 \\
Hub cannot suffice vs.\ can, constant & 2401--20 & 0.737 / 0.600 & +0.137 & [+0.106, +0.168] & $<$0.001 & +0.44 & 72/104 / 9/104 \\
Rare attribute cannot suffice, $1/\ln$ (H) & 2401--20 & 0.729 / 0.729 & 0.000 & [$-$0.028, +0.028] & 0.980 & $-$0.03 & 68/104 / 67/104 \\
$1/\ln(1+d)$ vs.\ constant, no cap & 2001--10 & 0.700 / 0.729 & $-$0.028 & [$-$0.068, +0.009] & 0.218 & 0.00 & 29/52 / 31/52 \\
$1/\ln$ vs.\ constant, both two-evidence & 2401--20 & 0.729 / 0.737 & $-$0.008 & [$-$0.020, +0.003] & 0.174 & $-$0.03 & 68/104 / 72/104 \\
Pairwise + cap 10 vs.\ rules + scorer & 2001--10 & 0.616 / 0.611 & +0.005 & [$-$0.039, +0.051] & 0.885 & $-$0.30 & 9/52 / 27/52 \\
\bottomrule
\end{tabular}}
\end{table*}

Table~\ref{tab:attribution} isolates the parts of the pairwise rule inside the
full pipeline, on fresh seeds. Removing hub links with a degree cap of 10 costs
\RfCapD{} F1 (95\% CI \RfCapLo{} to \RfCapHi{}, confirmatory, $p=\RfCapP$),
\RfCapHub{} on \texttt{hub\_heavy}, where recall of hub-hiding rings falls from
\RfCapHHNo{} to \RfCapHHCap{}. Keeping hub links but letting a single attribute
link a pair (a constant weight with $\theta_e=0.5$) costs \RfNdOneD{}, and
\RfNdOneHub{} on \texttt{hub\_heavy}: every co-user of the hub is linked and the
ring cannot be separated from them. Under the condition that a busy attribute
cannot suffice alone, the weight's shape does not matter: $1/\ln(1+d)$ and a
constant weight differ by \RfLogConstD{} (CI \RfLogConstLo{} to \RfLogConstHi{})
without a cap, and by \RfNdShapeD{} (CI \RfNdShapeLo{} to \RfNdShapeHi{}) when
both also bar rare attributes from linking alone. That last restriction
(the confirmatory test of the evidence condition) changes nothing overall
(\RfNdPrimD{}, CI \RfNdPrimLo{} to \RfNdPrimHi{}) but trades scenarios: it removes
almost all false rings on \texttt{drift} (\RfNdDriftFalseOne{} to
\RfNdDriftFalseTwo{} per run; F1 \RfNdDriftD{}) and loses adaptive rings
(\RfNdAdaptOne{} to \RfNdAdaptTwo{}; \texttt{adversarial} \RfNdAdvD{}), which
often link through one rarely shared device. Finally, with the same cap the
pairwise pipeline and the scored rules tie overall (\RfPairCapVsRulesD{}) but
trade sharply: the pipeline wins under benign noise (\RfPairCapVsRulesNoise{} on
\texttt{high\_noise}) and the rules win on hub hiding
(\RfPairCapVsRulesHub{} on \texttt{hub\_heavy}). The gain over the rules in
Table~\ref{tab:baselines} therefore has two sources: pairwise evidence with
several detectors, which helps under noise, and keeping hub links, which helps
under hub hiding.

Two exploratory follow-ups on further fresh seeds (protocol in
\texttt{docs/research/round3\_protocol.md}) sharpen this. First, the cost of a cap
does not come only from attributes that could link a pair alone. A cap of 27
removes exactly the attributes that, at $\theta_e=0.3$, can never link a pair
without a second piece of evidence, and still costs \RtCapTwentySevenD{} F1 (CI
\RtCapTwentySevenLo{} to \RtCapTwentySevenHi{}; cap 50: \RtCapFiftyD{}; cap 10:
\RtCapTenD{}). Hub-hiding rings found fall from \RtCapHHNo{} to
\RtCapHHTwentySeven{}. Busy attributes therefore help as \emph{supporting}
evidence, added to a transfer or a rarer attribute. Second, the property that
rare attributes link alone is what the adaptive rings need, but a step function
with that property is not a substitute for the discount: $h(d)=1$ for $d\le27$ and
0.2 above costs \RtStepD{} (CI \RtStepLo{} to \RtStepHi{}), mostly on
\texttt{hub\_heavy}, where the flat 0.2 lets busy attributes add up. A constant
weight that requires two pieces of evidence ($\theta_e=1.0$, the constant row of
Table~\ref{tab:existing}) differs by \RtConstTwoD{} (CI \RtConstTwoLo{} to \RtConstTwoHi{}),
with fewer false rings, but finds \RtConstTwoAdapt{} adaptive rings against
\RtLogAdapt{} for $1/\ln(1+d)$.

\subsection{RQ2: latency and scale}
% generated from results/e4_scale/results.csv (sha256 of LF text 7eb8cf45e1ba0a9a, 30 rows)
\begin{table*}[t]
\centering
\caption{Per-batch detection latency (median, ms; mean $\pm$ 95\% CI over 3 seeds)
and users examined per batch, incremental vs full recomputation, in-memory store. With three seeds
the $t$-interval can exceed the mean; the text uses ranges over seeds.}
\label{tab:scale}
\begin{tabular}{rccccc}
\toprule
& \multicolumn{2}{c}{Incremental} & \multicolumn{2}{c}{Full} & \\
Users & Region & ms & Region & ms & Speed-up \\
\midrule
1,000 & 976 & 211.0 $\pm$ 142.9 & 1015 & 159.7 $\pm$ 64.4 & 0.8$\times$ \\
5,000 & 1878 & 177.8 $\pm$ 46.7 & 5016 & 589.8 $\pm$ 67.3 & 3.3$\times$ \\
20,000 & 1324 & 54.6 $\pm$ 37.1 & 20010 & 2683.7 $\pm$ 1062.0 & 49.2$\times$ \\
50,000 & 1110 & 38.5 $\pm$ 38.2 & 50011 & 6181.6 $\pm$ 2576.4 & 160.4$\times$ \\
100,000 & 827 & 36.1 $\pm$ 44.2 & 100013 & 22500.5 $\pm$ 10278.5 & 623.0$\times$ \\
\bottomrule
\end{tabular}
\end{table*}

\begin{figure}[t]
\centering
\includegraphics[width=\columnwidth]{\figpath/fig2_latency_scaling.pdf}
\caption{Median per-batch detection latency as the graph grows (mean over 3
seeds; bars: per-seed minimum and maximum), incremental (dirty region) vs.\ full recomputation of the same
detectors, in-memory store.}
\label{fig:scale}
\end{figure}

Table~\ref{tab:scale} and Fig.~\ref{fig:scale} grow the user population of the
\texttt{scale\_sweep} scenario at a fixed event rate (100 events/s) and compare
incremental detection with recomputing the same detectors on the whole graph
for every batch. Full recomputation grows at least linearly with the graph
(faster between 50{,}000 and 100{,}000 users, where the 7.7\,GB test machine was
short of memory; we did not separate the two effects). The incremental region
stays between a few hundred and two thousand users, because the fan-out limits
stop the expansion at hubs; it even shrinks as the graph grows, since this
scenario scales the number of devices, cards and IPs with the users (so an
ordinary attribute is shared by fewer users) while hub sizes stay fixed. Its
latency therefore stays flat: at \ScaleMaxUsers{} users a batch takes \ScaleIncMs{}\,ms
incrementally against \ScaleFullMs{}\,ms in full. This comparison is a worst-case
bound, not a deployment comparison: nobody would recompute the whole graph on
every 50-event batch, the incremental figure (3 seeds) varies more than its
mean, and the 50{,}000 and 100{,}000-user full figures are confounded by memory
pressure. At
1{,}000 users the region already covers almost the whole graph, and the
incremental mode is slower than full recomputation because of the cost of
expanding the region. On the accuracy scenarios, incremental and full
recomputation flag exactly the same rings in every run (ablation
\emph{no incremental} in Table~\ref{tab:ablation}), so restricting detection to
the dirty region loses nothing there.

\textbf{Against a timer.} The realistic alternative is to run the same
detectors on the whole graph on a timer. On the accuracy scenarios (30 to 200
users, \TmRuns{} paired runs) a 1\,s timer matches the incremental engine in F1
(\TmOneF{} vs.\ \TmIncF{}, $p=\TmOnePval{}$) and time to detect
(\TmOneTTD{}\,s vs.\ \TmIncTTD{}\,s) at lower detection cost (\TmOneCost{}
vs.\ \TmIncCost{}\,ms per simulated second); a 5\,s timer costs less still
(\TmFiveCost{}) but detects later (\TmFiveTTD{}\,s), and a 30\,s timer misses
almost every ring, which ends before two passes confirm it (F1 \TmThirtyF{}).
On small graphs, incremental detection therefore buys nothing. The picture
reverses with size. In \texttt{scale\_sweep} with hubs that grow with the
population (15\% and 8\% of users, other hubs up to 10\%; \GhSeeds{} seeds), the
cost per simulated second is \GhIncCostA{}, \GhIncCostB{} and \GhIncCostC{}\,ms
for incremental detection at 1{,}000, 5{,}000 and 20{,}000 users, against
\GhTimerCostA{}, \GhTimerCostB{} and \GhTimerCostC{}\,ms for a 5\,s timer
(\GhFullCostC{}\,ms for recomputing every batch at 20{,}000). At 5{,}000 users
the two overlap across seeds; at 20{,}000 users the timer needs more than one
second of computation per simulated second in all three seeds, so it cannot
keep up, while the incremental region stays near \GhIncRegionC{} users even
though the hubs grow. On five fresh seeds at 10{,}000 users (revision
container, run alone) the 5\,s timer needs \RfScaleTenTimerLo{}--\RfScaleTenTimerHi{}\,ms
per simulated second, more than one second in \RfScaleTenTimerOver{} of 5 seeds;
the incremental engine is cheaper in \RfScaleTenWins{} of 5 (median
\RfScaleTenIncMed{}\,ms), but in one seed it costs \RfScaleTenIncHi{}\,ms, more
than the timer. The incremental
cost is not monotone in size (\GhIncCostB{} at 5{,}000, \GhIncCostC{} at
20{,}000); a likely cause, which we did not isolate, is that the region is
bounded by the fan-out limits while the number of touched users per batch
falls as the graph grows at a fixed event rate. We did not run 50{,}000 users
with a timer. The static rules run on
the same 5\,s timer, which is why their time to detect is similar to ours on the
accuracy scenarios (\EPropTTD{}\,s vs.\ \ECapTenTTD{}\,s).

\subsection{RQ3: robustness to evasion}
\begin{figure}[t]
\centering
\includegraphics[width=\columnwidth]{\figpath/fig4_hub_size.pdf}
\caption{A shared-IP ring hiding behind an office hub of growing size
(\texttt{hub\_heavy}, 5 test seeds per size). Left: recall of the hub-hiding
rings; right: ring F1 (mean, 95\% CI).}
\label{fig:hubsize}
\end{figure}
\begin{figure*}[t]
\centering
\includegraphics[width=\textwidth]{\figpath/fig5_evasion.pdf}
\caption{Ring recall by evasion strategy, pooled over all runs of
Table~\ref{tab:baselines} (static rules without a scorer; Table~\ref{tab:existing}
gives the recall by strategy for every method with its own scorer).}
\label{fig:evasion}
\end{figure*}

\textbf{Hub size.} In Fig.~\ref{fig:hubsize} the office hub grows from
\HubSizeMin{} to \HubSizeMax{} accounts. Recall of the hub-hiding rings does
not fall with hub size for any method (proposed \HubSizePropHidden{}, static
cap 10 \HubSizeCapTenHidden{} over all sizes). The static rules do not find
these rings through the shared hub: with the 3-cycle rule switched off, cap 10
finds \HubSizeCapTenNoCycleHidden{}, because the rings also move money among
themselves in short cycles. What separates the methods is precision, and so
F1 (\HubSizePropF{} vs.\ \HubSizeCapTenF{} for the unscored rules;
\RfScRulesHubSizeF{} for the rules with their own scorer).

\textbf{Hubs above the fan-out limit.} Fig.~\ref{fig:hubsize} stops at 200
accounts, below the 250-user first-hop limit. We therefore resized the office
hub to 500, 2,000 and 10,000 accounts (5 seeds each) and compared the shipped
engine with full recomputation and with a ``bound'' rule that does cross an
over-limit attribute, towards those co-members whose attribute weight with the
touched user could reach the edge threshold. At every size the bound rule
flagged the same rings as the 250-user limit, with identical precision, recall
and region size (\LhBoundSame{} runs): it never added a user, so no pair in these
streams shared two over-limit attributes. Full recomputation does not give
identical runs. Tracing them showed two causes, neither of them the limit: a
ring none of whose members is touched enters a region only partly, and greedy
peeling of a connected component of several hundred users can return the ring
on the whole graph but not on the region's part of it. In total it found no
more hub-hiding rings (\LhLogFullHHA{} vs.\
\LhLogFanHHA{} at 500, \LhLogFullHHB{} vs.\ \LhLogFanHHB{} at 2,000) and its F1
was not higher (\LhLogFullFA{} vs.\ \LhLogFanFA{}, \LhLogFullFB{} vs.\
\LhLogFanFB{}), while the incremental region stayed at a median of
\LhRegionInc{} users against up to \LhRegionFullMax{}. At 10,000 accounts one
full-recomputation run did not finish within 3.5\,h on the test laptop, so
there is no full-recomputation reference at that size. A constant weight (edge threshold 1.0, as selected in RQ5) found fewer
hub-hiding rings incrementally than with full recomputation at 500 accounts
(\LhConstFanHHA{} vs.\ \LhConstFullHHA{}; F1 \LhConstFanFA{} vs.\
\LhConstFullFA{}, lower in 3 of 5 runs; five seeds cannot reach significance),
by the same two mechanisms; a single busy attribute cannot reach this threshold
on its own. One full run of it took about \LhConstFullWallB{} minutes at 2,000
accounts. So, for hubs up to 10,000 accounts, the 250-user limit itself hid no
pair for either weight, but ring-level equality with full recomputation does
not hold at these sizes. The static rules stay behind (F1 \LhCapTenFA{},
\LhCapTenFB{}, \LhCapTenFC{} without a scorer; \RfScRulesLhA{},
\RfScRulesLhB{}, \RfScRulesLhC{} with their own scorer).

\textbf{Evasion strategies.} Fig.~\ref{fig:evasion} breaks recall down by the
strategy each planted or random ring uses. The proposed method's advantage in
recall comes mostly from rings without evasion (\EvPropNone{} vs.\
\EvCapTenNone{} for cap 10) and rings with decoy edges (\EvPropDecoyEdges{}
vs.\ \EvCapTenDecoyEdges{}). On hub-hiding rings the two are close
(\EvPropHubHiding{} vs.\ \EvCapTenHubHiding{}), and identity rotation
(\EvPropIdentityRotation{}) and low-and-slow activity (\EvPropLowAndSlow{})
defeat every method we tested, presumably because rotation removes the shared
attributes the detectors rely on and slow rings accumulate little evidence
inside one window.

\textbf{Adaptive attacker.} With feedback switched on, the attacker is told
which of its rings were alerted and shifts its strategy mix. Over
\AdaptSeeds{} test seeds it switched strategy \AdaptOnSwitches{} times and
launched \AdaptOnRings{} adaptive rings, of which \AdaptOnFound{} were
detected; without feedback, \AdaptOffFound{} of \AdaptOffRings{}. We saw no
change in the detection rate (difference 0.00, 95\% interval \RfAdaptLo{} to
\RfAdaptHi{}), an interval too wide to rule out an effect, and the attacker's
strategies are limited to the four we implemented.

\subsection{RQ4: sensitivity}
\begin{figure*}[t]
\centering
\includegraphics[width=\textwidth]{\figpath/fig6_sweeps.pdf}
\caption{Ring F1 of the proposed method and of the static rules (cap 10,
without a scorer) when the event rate (\texttt{adversarial}), the planted ring
size (\texttt{baseline}) and the view window (\texttt{drift}, proposed only)
change; 5 test seeds per point. With their own scorer the rules reach F1
\RfScRulesRateLo{} at the lowest and \RfScRulesRateHi{} at the highest event
rate, and \RfScRulesSizeLo{} at the smallest and \RfScRulesSizeHi{} at the
largest ring size.}
\label{fig:sweeps}
\end{figure*}
Raising the event rate from \SwRateLo{} to \SwRateHi{} events/s lowers the
proposed method's F1 from \SwRateLoF{} to \SwRateHiF{} and raises its per-batch
latency from \SwRateLoLat{} to \SwRateHiLat{}\,ms, since each batch touches more
of the graph (Fig.~\ref{fig:sweeps}). Larger planted rings are harder (F1
\SwRingSizeLoF{} at size \SwRingSizeLo{}, \SwRingSizeHiF{} at size
\SwRingSizeHi{}); a likely cause, which we did not isolate, is the 0.75 edge
density the peeling requires. The view window between
\SwWindowLo{} and \SwWindowHi{}\,s changes F1 little (\SwWindowLoF{} vs.\
\SwWindowHiF{}).

\subsection{RQ5: ablations}
% generated from results/e2_ablations/results.csv (sha256 of LF text ac950ecba52299df, 180 rows)
\begin{table*}[t]
\centering
\caption{Ablations on adversarial, hub\_heavy, high\_noise, baseline
($\times$ 5 test seeds). $\Delta$F1: paired mean difference to the full system;
$d_z$: paired effect size of the full system over the variant; $p$: Wilcoxon signed-rank,
uncorrected and Holm-corrected over the 7 tested comparisons.}
\label{tab:ablation}
\begin{tabular}{lcccccccc}
\toprule
Variant & Precision & Recall & F1 & $\Delta$F1 & $d_z$ & $p$ & $p_{\mathrm{Holm}}$ & Latency (ms) \\
\midrule
Full system & 0.75 $\pm$ 0.12 & 0.73 $\pm$ 0.12 & 0.72 $\pm$ 0.11 & -- & -- & -- & -- & 19.2 \\
w/o hub weighting & 0.65 $\pm$ 0.15 & 0.58 $\pm$ 0.14 & 0.60 $\pm$ 0.14 & -0.12 & 0.42 & 0.161 & 0.462 & 24.0 \\
w/o time decay & 0.80 $\pm$ 0.11 & 0.72 $\pm$ 0.10 & 0.75 $\pm$ 0.09 & +0.02 & -0.18 & 0.277 & 0.462 & 17.6 \\
w/o incremental scheduling & 0.75 $\pm$ 0.12 & 0.73 $\pm$ 0.12 & 0.72 $\pm$ 0.11 & 0.00 & -- & -- & -- & 18.4 \\
w/o risk model & 0.28 $\pm$ 0.09 & 0.82 $\pm$ 0.09 & 0.39 $\pm$ 0.10 & -0.33 & 2.05 & $<$0.001 & $<$0.001 & 15.1 \\
w/o hysteresis & 0.72 $\pm$ 0.12 & 0.74 $\pm$ 0.10 & 0.70 $\pm$ 0.10 & -0.02 & 0.37 & 0.154 & 0.462 & 22.7 \\
w/o hub-aware detector & 0.98 $\pm$ 0.05 & 0.25 $\pm$ 0.09 & 0.36 $\pm$ 0.12 & -0.36 & 1.02 & 0.001 & 0.006 & 5.3 \\
w/o cycle detector & 0.75 $\pm$ 0.14 & 0.60 $\pm$ 0.11 & 0.64 $\pm$ 0.11 & -0.09 & 1.02 & 0.002 & 0.009 & 15.4 \\
w/o smurfing detector & 0.72 $\pm$ 0.14 & 0.64 $\pm$ 0.15 & 0.65 $\pm$ 0.14 & -0.07 & 0.93 & 0.003 & 0.013 & 19.9 \\
\bottomrule
\end{tabular}
\end{table*}

Table~\ref{tab:ablation} removes one component at a time (\AblPairs{} paired
runs). The calibrated scorer contributes most: without it F1 falls from
\AblFullF{} to \AblNoScoringF{} because precision drops to \AblNoScoringP{}.
Without the hub-aware detector, F1 falls to \AblNoHubAwareF{} and to
\AblNoHubAwareFHubHeavy{} on \texttt{hub\_heavy}, but precision rises to
\AblNoHubAwareP{} at recall \AblNoHubAwareR{}: most false candidates come from
the hub-aware detector, and the scorer removes most of them. This suggests a
separate alert threshold for rings the hub-aware detector found; selected on
the validation seeds over a $6\times6$ grid, the best pair was the shipped
single threshold (0.8, 0.8), so the refinement gains nothing here. Removing the cycle or the
smurfing detector costs \AblNoCycleF{} and \AblNoSmurfingF{}. Replacing the hub
weight by a constant at the same edge threshold collapses \texttt{hub\_heavy}
(\AblFullFHubHeavy{} to \AblNoHubWeightFHubHeavy{}) but helps
\texttt{high\_noise} (\AblFullFHighNoise{} to \AblNoHubWeightFHighNoise{}). The
effect depends on the scenario (heterogeneity $p=\RfAblHubWHet$), so the pooled
test ($p=\AblNoHubWeightPval{}$) is not informative either way. That ablation keeps the edge threshold tuned for
$1/\ln(1+d)$, which is unfair to the constant; the next paragraph gives each
weight its own threshold.

\textbf{Weight functions.} We compared five weights $h(d)$, each with its edge
threshold selected on the validation seeds and the alert threshold fixed at
\TunedAlert{}: the shipped $1/\ln(1+d)$ ($\theta_e=\WShipEdge{}$), $1/d$
(\WInvEdge{}), $1/\sqrt{d}$ (\WSqrtEdge{}), a soft cap
$1/(1+e^{(d-10)/2})$ (\WSoftEdge{}) and a constant (\WConstEdge{}); for $1/d$ and the soft cap the selected value lies on the edge of this grid
(\RfGridNote{}); for the constant, a wider grid in Section~\ref{sec:attribution}
selects 1.0 again as an interior optimum (1.2 is worse).
On \texttt{hub\_heavy} with \WbSeeds{} paired seeds, the shipped weight reaches
F1 \WbShipF{} and recall \WbShipHH{} of the rings hidden behind hubs; the
constant reaches \WbConstF{} and \WbConstHH{} ($p=\WbConstPval{}$), $1/\sqrt{d}$
\WbSqrtF{}, while $1/d$ (\WbInvF{}) and the soft cap (\WbSoftF{}, hidden-ring
recall \WbSoftHH{} at precision \WbSoftP{}) are worse ($p$ \WbSoftPval{}). Over
the five scenarios the picture is the same: \WcShipF{} for the shipped weight,
\WcConstF{} for the constant and \WcSqrtF{} for $1/\sqrt{d}$, and \WcInvF{} for
$1/d$ and \WcSoftF{} for the soft cap. With the widened grid, $1/d$ at
$\theta_e=0.05$ reaches \RfInvFiveF{} over the five scenarios (\RfInvFiveD{}
against the shipped weight, 95\% CI \RfInvFiveLo{} to \RfInvFiveHi{}) and
\RfInvHubF{} on \texttt{hub\_heavy} (\RfInvHubD{}, CI \RfInvHubLo{} to
\RfInvHubHi{}), so most of its earlier deficit came from the grid edge.
$1/\sqrt{d}$ is equivalent to the shipped
weight within 0.05 F1 (90\% interval \RfWSqrtLo{} to \RfWSqrtHi{}); the constant
is not significantly different but not shown equivalent (\RfWConstLo{} to
\RfWConstHi{}). Switching hub links off loses rings that share a hub
(Section~\ref{sec:attribution}); a strong discount costs little once the edge
threshold is tuned for it, and the soft cap, which nearly switches busy
attributes off, stays behind.
Hysteresis makes no difference beyond 0.05 F1 (90\% interval \RfAblHystLo{} to
\RfAblHystHi{}); time decay makes no significant difference but is not shown
equivalent (\RfAblDecayLo{} to \RfAblDecayHi{}), and without decay F1 is even
slightly higher (\AblNoTimeDecayF{}). The accuracy scenarios last
60--90\,s, short against the 300\,s half-life, so decay has little to act on
here; its value on long streams is untested.

\subsection{Public data}
\label{sec:public}
\input{sections/public.tex}

\section{Deployment, Calibration and Explanation Studies}
\label{sec:phase2}
The experiments in this section were added after the main evaluation. Each has
a protocol committed to the repository before its first result
(\texttt{docs/research/}), and every amendment is recorded there with its
reason.

\subsection{RQ6: latency on the deployed stack}
The latencies of Section~\ref{sec:results} come from the in-memory store. We
therefore replayed the paced scenarios of RQ1 and the users sweep through
the deployed path: Redis Streams into the detection engine, with the graph in
Neo4j~5.26. All \EightRuns{} runs produced the same sequence of ring updates and
the same flagged rings as the in-memory engine (\EightSame{} of \EightRuns{}),
so the store and the bus change timing, never results. Paced at simulated real
time, a flag followed the event that completed its batch by a median of
\EightFlagLo{}--\EightFlagHi{}\,ms across the five scenarios. Each detection
reads its region through Cypher round trips and took \EightSlowLo{} to
\EightSlowHi{} times as long as on the in-memory store for the same batches.
When events arrive faster than the store can take them, Neo4j sets the ceiling:
\EightTputLo{}--\EightTputHi{} events per second on the test laptop, with events
queueing in Redis and flags arriving \EightBurstLo{}--\EightBurstHi{}\,s after
their event. Varying the consumer's read size and blocking timeout, after
HYADES~\cite{bora2026hyades}, changed nothing measurable at these rates. Neo4j
saturates before the consumer does.

\subsection{RQ7: ring-level calibration and early warning}
The scorer's isotonic calibration was fitted on snapshot rows, but analysts read
scores of ring-level alerts. On \TenUnits{} ring-level units from the test seeds
its expected calibration error is \TenEce{}, and rings scored in the 0.8--0.9
and 0.9--1.0 bins were true rings \TenBinEight\% and \TenBinNine\% of the time.
A calibrator refitted on ring-level units of separate seeds, in the manner
compared by Azbergen et al.~\cite{azbergen2026calibration} and Ding et
al.~\cite{ding2021temperature}, lowers the Brier score from \TenBrier{} to
\TenBrierRefit{}. Ranking holds up better than calibration: precision at the top
5\% of units is \TenPAtFive{}. Early warning, measured without look-ahead, is
rare. Of \ElevenEligible{} true rings with a forming phase, \ElevenWarned{}
(\ElevenShare\%) were flagged before activation, a median of \ElevenLead{}\,s
early, and \ElevenPrecision\% of first early warnings pointed at a real ring.

\subsection{RQ8: explanation quality}
We measured the explanations the product shows on the \TwelveAlerts{} alerts of
the test streams. The unit is an alert (a ring at its peak calibrated score of
at least \TunedAlert{}); a model counterfactual is a proposed change of one
actionable feature, and an alert can have none, one or several.
\TwelveCfAlerts{} of the \TwelveAlerts{} alerts have at least one, \TwelveCfs{}
in all. Writing the unit test for counterfactual validity exposed a defect:
counterfactuals aimed at the raw score of the calibrator's breakpoint, but the
calibrator interpolates between breakpoints, so the proposed change left the
calibrated score above the threshold. As shipped, \TwelveValidAsShipped{} of
the \TwelveCfs{} counterfactuals flipped the decision; aimed at the
interpolated crossing, \TwelveValidCorrected{} of \TwelveCfs{} do. The
correction was registered as a protocol amendment before the experiment ran.
Any system that derives counterfactual targets from an isotonic calibrator can
make the same mistake.

Coverage is the larger limitation: the other \TwelveNoCfShare\% of alerts have
no single-feature counterfactual over the actionable features, because typical
transfer size dominates their score. Under 5\% noise in non-actionable
features, after Artelt et al.~\cite{artelt2021robustness}, the contribution
vector moves by a median of \TwelveDriftFive{} (relative $L_1$) and the top
three features never change for the median alert; for a linear model this is
close to guaranteed, so it mainly checks the implementation. A scorer retrained
on disjoint seeds, after Movin et al.~\cite{movin2025acdc}, ranks the
contributions similarly (median Spearman \TwelveSpearman{}) and still raises
\TwelveStillAlert\% of the alerts, but it proposes the same first
counterfactual for only \TwelveCfSame\% of them, so a counterfactual describes
this model, not the ring. Against simulator ground truth, the members shown in
the evidence subgraph of a correct alert are the true ring's (median precision
\TwelveFidPrec{}, coverage \TwelveFidCov{}). Because that subgraph is built from
the flagged ring's members, this checks that the evidence shows the right
accounts; it does not test whether the evidence explains the score. A
deletion test would: remove the shared-attribute links or the internal
transfers that the evidence lists, recompute the ring's features and score, and
compare the drop with removing as many links of the same members that the
evidence does not list. We have not run it.

\subsection{RQ9: a learned baseline and evasion strength}
To test whether the rule stage is needed, we trained gradient boosting on the
same 19 ring features over rule-free candidate groups (Louvain communities
without a density filter). Training, calibration and threshold selection used
the shipped scorer's protocol. On the test streams of RQ1 it reaches F1
\ThirteenGbF{} against \ThirteenPropF{} (paired Wilcoxon $p\ThirteenP$,
rank-biserial \ThirteenRb{}). Precision is not the problem: it is \ThirteenGbP{}
against \ThirteenPropP{}, with a third of the alerts. Recall is:
\ThirteenGbR{} against \ThirteenPropR{}. The gap is largest where hubs and
evasion matter (\texttt{hub\_heavy}: \ThirteenGbFHub{} against
\ThirteenPropFHub{}; \texttt{adversarial}: \ThirteenGbFAdv{} against
\ThirteenPropFAdv{}).

Turning up camouflage strength, after Liu et al.~\cite{liu2025robust} and Hu et
al.~\cite{hu2025hhgfd}, separates the methods further. With up to eight decoy
edges per ring action, the proposed method's recall of the decoy rings stays
between \FourteenDecoyPropLo{} and \FourteenDecoyPropHi{}, while Louvain falls
to \FourteenDecoyLouvainMax{} and the learned baseline to
\FourteenDecoyGbMax{}. Slow rings remain the weakness found in RQ3: the proposed
method recalls at most \FourteenSlowPropHi{} of them, and static Cypher with a
degree cap of ten recovers more (\FourteenSlowCapLo{}--\FourteenSlowCapHi{}),
though at a precision near 0.04. In the 60\,s scenario, strong slow and
identity-rotation settings produce few rings, so those two sweeps have little
statistical power. We report them as indicative.


\section{Threats to Validity}
\label{sec:threats}
\textbf{Simulator realism.} All main results come from our simulator. Hub
sizes, ring behaviour and evasion strategies follow our configuration, not
observed fraud. A detector tuned to this generator may do worse on real
traffic, and the absolute numbers should not be read as field performance.

\textbf{Synthetic labels.} Ground truth is the set of rings the simulator
planted. Benign groups that happen to look like rings count as false
positives, and planted rings that never become visible in the stream still
count against recall; both directions affect every method equally.

\textbf{Public data.} The public datasets we could obtain carry no device,
card or IP links, so they cannot test hub-aware weighting at all; on them only
the transaction detectors run (Section~\ref{sec:public}). Only AMLworld has
group labels, and we could replay only a four-day slice with a one-hour window.
Elliptic and the AMLSim samples turned out to be outside what the transaction
rules can detect, for structural reasons we describe. A real-world test of
hub-aware weighting would need data with shared-attribute links, which we did
not find in public form.

\textbf{Tuning.} Thresholds were selected on validation seeds of the same
generator; the test seeds are independent draws but from the same
distribution. In the weight-function comparison, three of the five selected
edge thresholds lie on the edge of the grid, so those functions might do
somewhat better with a wider grid.

\textbf{Scale.} The largest graph we measured has \ScaleMaxUsers{} users. A
million users was not feasible on the 7.7\,GB test machine: at 100{,}000 users
less than 0.3\,GB of memory was free and one full recomputation already took
over 20\,s per batch. In the scale scenario, attribute counts grow with the
number of users while hub sizes stay fixed; graphs whose hubs grow with the
population would give larger dirty regions. The 250-user first-hop limit was
tested with hubs of up to 10,000 accounts, five seeds per size, on one scenario;
its harmlessness rests on the hub-aware weight keeping a single busy attribute
below the edge threshold, and would not hold for a weight or threshold that
does not.

\textbf{Latency.} The latencies of Section~\ref{sec:results} come from the
in-memory store on a laptop. Section~\ref{sec:phase2} repeats the measurement on
Neo4j and Redis Streams. The detected rings are identical, but detection is
several times slower and throughput is bounded by the store. Both figures
describe one laptop running Docker, not a server.

\textbf{Baselines.} The scored baselines of Table~\ref{tab:baselines} use the
shipped scorer, which was trained on candidates of the proposed pipeline; a
scorer trained on each baseline's own candidates might do better for it. The
learned baseline of Section~\ref{sec:phase2} is a
tabular model over Louvain groups, so its weakness under decoy edges is partly
the grouping's. We did not run a graph neural network. A GraphSAGE node
classifier with a fixed rule for turning node scores into rings, as used by
Badjie et al.~\cite{badjie2026ensembles} and Kumar et al.~\cite{kumar2026tagat},
is the obvious next comparison, together with node-level AUC on public
benchmarks. Our node-level numbers are not comparable to those benchmarks.

\textbf{Calibration.} The shipped scorer is miscalibrated at the ring level
(Section~\ref{sec:phase2}). Its scores rank alerts well but should not be read
as probabilities until a ring-level calibrator replaces it.

\section{Conclusion}
Keeping links through busy shared attributes, instead of capping them, lets a
streaming detector reach a recall on rings that share busy infrastructure that
no cap reaches. The shape of the discount matters little: a constant weight
with its own threshold did as well as $1/\ln(1+d)$, while strong discounts and
soft caps lost the rings hidden behind hubs. On its own the weighting is not
precise, and a calibrated scorer over ring features supplies most of the
precision. Given the same scorer, the
degree-capped rules reach F1 \ECapTenSF{} against \EPropF{} for the full
pipeline, so the scorer accounts for most of the gain over the original rules,
and at the best operating point on shared-attribute rings the weight and the
cap are not distinguishable on F1. Restricting detection to the region a
micro-batch touches gave the same rings as full recomputation in every
accuracy run, and kept per-batch latency flat as the graph grew to
\ScaleMaxUsers{} users while full recomputation grew past 20\,s. Rings that
rotate identities or act slowly evade all methods we tested, and all of our
positive results are on synthetic data. On Neo4j and Redis Streams the detector flags exactly the same rings, at the
cost of slower detection and a throughput bounded by the store. A learned model
over rule-free groups is more precise but finds half as many rings, and it
collapses under decoy edges, which the hub-aware rules withstand. The
explanations are stable and their evidence is faithful. Measuring them also
found and fixed a calibration defect that had made every counterfactual
invalid. The next steps are data with real shared-attribute links, a ring-level
calibrator, a GraphSAGE baseline, multi-feature counterfactuals, and longer runs
for slow and identity-rotating rings.

Code, scenarios, models and the scripts that regenerate every table and figure
are in the project repository (\texttt{REPRODUCE.md}).
